\documentclass[10pt,a4paper]{amsart}
\usepackage{amsmath,amssymb,amsthm,graphicx}
\usepackage{verbatim}
\newenvironment{reference}{\comment}{\endcomment}
%\newenvironment{reference}{}{}
\newenvironment{slogan}{\comment}{\endcomment}
\newenvironment{history}{\comment}{\endcomment}

% For commutative diagrams we use Xy-pic
\usepackage[all]{xy}

% We use 2cell for 2-commutative diagrams.
\xyoption{2cell}
% Theorem environments.
%
\theoremstyle{plain}
\newtheorem{theorem}[subsection]{Theorem}
\newtheorem{proposition}[subsection]{Proposition}
\newtheorem{lemma}[subsection]{Lemma}

\theoremstyle{definition}
\newtheorem{definition}[subsection]{Definition}
\newtheorem{example}[subsection]{Example}
\newtheorem{exercise}[subsection]{Exercise}
\newtheorem{situation}[subsection]{Situation}

\theoremstyle{remark}
\newtheorem{remark}[subsection]{Remark}
\newtheorem{remarks}[subsection]{Remarks}

\numberwithin{equation}{subsection}

% Macros
%
\def\lim{\mathop{\mathrm{lim}}\nolimits}
\def\colim{\mathop{\mathrm{colim}}\nolimits}
\def\Spec{\mathop{\mathrm{Spec}}}
\def\Hom{\mathop{\mathrm{Hom}}\nolimits}
\def\Ext{\mathop{\mathrm{Ext}}\nolimits}
\def\SheafHom{\mathop{\mathcal{H}\!\mathit{om}}\nolimits}
\def\SheafExt{\mathop{\mathcal{E}\!\mathit{xt}}\nolimits}
\def\Sch{\mathit{Sch}}
\def\Mor{\mathop{\mathrm{Mor}}\nolimits}
\def\Ob{\mathop{\mathrm{Ob}}\nolimits}
\def\Sh{\mathop{\mathit{Sh}}\nolimits}
\def\NL{\mathop{N\!L}\nolimits}
\def\CH{\mathop{\mathrm{CH}}\nolimits}
\def\proetale{{pro\text{-}\acute{e}tale}}
\def\etale{{\acute{e}tale}}
\def\QCoh{\mathit{QCoh}}
\def\Ker{\mathop{\mathrm{Ker}}}
\def\Im{\mathop{\mathrm{Im}}}
\def\Coker{\mathop{\mathrm{Coker}}}
\def\Coim{\mathop{\mathrm{Coim}}}

% Boxtimes
%
\DeclareMathSymbol{\boxtimes}{\mathbin}{AMSa}{"02}

%
% Macros for moduli stacks/spaces
%
\def\QCohstack{\mathcal{QC}\!\mathit{oh}}
\def\Cohstack{\mathcal{C}\!\mathit{oh}}
\def\Spacesstack{\mathcal{S}\!\mathit{paces}}
\def\Quotfunctor{\mathrm{Quot}}
\def\Hilbfunctor{\mathrm{Hilb}}
\def\Curvesstack{\mathcal{C}\!\mathit{urves}}
\def\Polarizedstack{\mathcal{P}\!\mathit{olarized}}
\def\Complexesstack{\mathcal{C}\!\mathit{omplexes}}
% \Pic is the operator that assigns to X its picard group, usage \Pic(X)
% \Picardstack_{X/B} denotes the Picard stack of X over B
% \Picardfunctor_{X/B} denotes the Picard functor of X over B
\def\Pic{\mathop{\mathrm{Pic}}\nolimits}
\def\Picardstack{\mathcal{P}\!\mathit{ic}}
\def\Picardfunctor{\mathrm{Pic}}
\def\Deformationcategory{\mathcal{D}\!\mathit{ef}}
\usepackage{cleveref}
\crefname{originalreference}{original reference}{original references}
\Crefname{originalreference}{Original reference}{Original references}
\makeatletter
\newlabel{remark-relative-dualizing-complex}{{0B6S}{source}{}{sourcepointer-1}{}}
\newlabel{remark-relative-dualizing-complex@cref}{{[originalreference][0][]0B6S}{[1][0][]source}{}{}{}}
\makeatother
\crefname{originallemma}{original lemma}{original lemmas}
\Crefname{originallemma}{Original lemma}{Original lemmas}
\makeatletter
\newlabel{lemma-twisted-inverse-image}{{0A9E}{source}{}{sourcepointer-2}{}}
\newlabel{lemma-twisted-inverse-image@cref}{{[originallemma][0][]0A9E}{[1][0][]source}{}{}{}}
\makeatother
\crefname{originallemma}{original lemma}{original lemmas}
\Crefname{originallemma}{Original lemma}{Original lemmas}
\makeatletter
\newlabel{lemma-compare-with-pullback-flat-proper}{{0E4K}{source}{}{sourcepointer-3}{}}
\newlabel{lemma-compare-with-pullback-flat-proper@cref}{{[originallemma][0][]0E4K}{[1][0][]source}{}{}{}}
\makeatother
\crefname{originallemma}{original lemma}{original lemmas}
\Crefname{originallemma}{Original lemma}{Original lemmas}
\makeatletter
\newlabel{lemma-proper-flat-base-change}{{0AAB}{source}{}{sourcepointer-4}{}}
\newlabel{lemma-proper-flat-base-change@cref}{{[originallemma][0][]0AAB}{[1][0][]source}{}{}{}}
\makeatother
\title{Rigidity: the relative dualizing complex inverted along the diagonal gives back the structure sheaf and the complex}
\author{The Stacks Project Authors}
\date{stacks-project:duality.tex:local-snapshot-20261012; GNU FDL 1.2 or later}
\begin{document}
\maketitle
\noindent\textit{Editorial note.} All delivered mathematical text is the author's own, copied from the pinned Stacks Project chapter file duality.tex; no statement or proof step was rewritten, reordered or removed.\par\medskip
\noindent\textit{Editorial note.} Preamble provenance: the pinned chapter file opens with an input of the shared preamble and carries no preamble text of its own, so the project's single Stacks-wide preamble staged at sources/stacks-duality/preamble.tex is excerpted verbatim under the recipe fixed by LANE\_BRIEF 22.1, as three spans, lines 8-13 (verbatim package and the reference/slogan/history environments), lines 14-18 (Xy-pic setup) and lines 154-217 (the native theorem declarations and mathematical macros). Nothing was added, removed or redefined; the wrapper class and the amsmath/amssymb/amsthm/graphicx base packages are supplied by the assembler.\par\medskip
\noindent\textit{Editorial note.} The counted statement and proof contain no cross reference, no citation and no included file, so the delivered body is self contained.\par\medskip
\noindent\textit{Editorial note.} Canonical location of the counted proposition: Stacks tag 0E2P (sources/stacks-duality/pinned-original-source.tex line 3159).\par\medskip
\noindent\textit{Allowed standard prerequisites (editorial).}\begin{itemize}
\item The right adjoints of pushforward and their composition along the diagonal. Same chapter, Lemma labelled lemma-twisted-inverse-image, cited at proof lines 3225-3235. Define a, b, c and the identity id = c o b.
\item The comparison with pullback for a flat proper morphism. Same chapter, Lemma labelled lemma-compare-with-pullback-flat-proper, cited at proof lines 3237-3240. Gives the second term of the chain.
\item Base change of the adjoint along the projection. Same chapter, Lemma labelled lemma-proper-flat-base-change, cited at proof lines 3240-3244. Gives the third term.
\item The relative dualizing complex as the adjoint applied to the structure sheaf. Same chapter, Remark labelled remark-relative-dualizing-complex, invoked in the statement, used at proof lines 3244-3248. Identifies a(O\_Y) with omega.
\end{itemize}\medskip
\ifdefined\hypertarget\hypertarget{sourcepointer-1}{}\fi
\noindent\textit{Source reference.} Original reference 0B6S: sources/stacks-duality/pinned-original-source.tex, source lines 2945--2945. Same-chapter pointer into the author environment remark at source line 2945; invoked at proof line None of the counted proof. It is retained only as a navigation pointer; every load-bearing step is inside the delivered span.\par\medskip
\ifdefined\hypertarget\hypertarget{sourcepointer-2}{}\fi
\noindent\textit{Source reference.} Original lemma 0A9E: sources/stacks-duality/pinned-original-source.tex, source lines 653--653. Same-chapter pointer into the author environment lemma at source line 653; invoked at proof line 3186 of the counted proof. It is retained only as a navigation pointer; every load-bearing step is inside the delivered span.\par\medskip
\ifdefined\hypertarget\hypertarget{sourcepointer-3}{}\fi
\noindent\textit{Source reference.} Original lemma 0E4K: sources/stacks-duality/pinned-original-source.tex, source lines 2892--2892. Same-chapter pointer into the author environment lemma at source line 2892; invoked at proof line 3217 of the counted proof. It is retained only as a navigation pointer; every load-bearing step is inside the delivered span.\par\medskip
\ifdefined\hypertarget\hypertarget{sourcepointer-4}{}\fi
\noindent\textit{Source reference.} Original lemma 0AAB: sources/stacks-duality/pinned-original-source.tex, source lines 2929--2929. Same-chapter pointer into the author environment lemma at source line 2929; invoked at proof line 3218 of the counted proof. It is retained only as a navigation pointer; every load-bearing step is inside the delivered span.\par\medskip
\begin{lemma}[Rigidity]
\label{lemma-van-den-bergh}
Let $Y$ be a quasi-compact and quasi-separated scheme.
Let $f : X \to Y$ be a proper, flat morphism of finite presentation
with relative dualizing complex $\omega_{X/Y}^\bullet$
(Remark \ref{remark-relative-dualizing-complex}).
There is a canonical isomorphism
\begin{equation}
\label{equation-pre-rigid}
\mathcal{O}_X =
c(L\text{pr}_1^*\omega_{X/Y}^\bullet) =
c(L\text{pr}_2^*\omega_{X/Y}^\bullet)
\end{equation}
and a canonical isomorphism
\begin{equation}
\label{equation-rigid}
\omega_{X/Y}^\bullet =
c\left(L\text{pr}_1^*\omega_{X/Y}^\bullet
\otimes_{\mathcal{O}_{X \times_Y X}}^\mathbf{L}
L\text{pr}_2^*\omega_{X/Y}^\bullet\right)
\end{equation}
where $c$ is the right adjoint of
Lemma \ref{lemma-twisted-inverse-image}
for the diagonal $\Delta : X \to X \times_Y X$.
\end{lemma}

\begin{proof}
Let $a$ be the right adjoint to $Rf_*$ as in
Lemma \ref{lemma-twisted-inverse-image}.
Consider the cartesian square
$$
\xymatrix{
X \times_Y X \ar[r]_q \ar[d]_p & X \ar[d]_f \\
X \ar[r]^f & Y
}
$$
Let $b$ be the right adjoint for $p$
as in Lemma \ref{lemma-twisted-inverse-image}. Then
\begin{align*}
\omega_{X/Y}^\bullet
& =
c(b(\omega_{X/Y}^\bullet)) \\
& =
c(Lp^*\omega_{X/Y}^\bullet
\otimes_{\mathcal{O}_{X \times_Y X}}^\mathbf{L} b(\mathcal{O}_X)) \\
& =
c(Lp^*\omega_{X/Y}^\bullet
\otimes_{\mathcal{O}_{X \times_Y X}}^\mathbf{L}
Lq^*a(\mathcal{O}_Y)) \\
& =
c(Lp^*\omega_{X/Y}^\bullet
\otimes_{\mathcal{O}_{X \times_Y X}}^\mathbf{L}
Lq^*\omega_{X/Y}^\bullet)
\end{align*}
as in (\ref{equation-rigid}). Explanation as follows:
\begin{enumerate}
\item The first equality holds as $\text{id} = c \circ b$ because
$\text{id}_X = p \circ \Delta$.
\item The second equality holds by
Lemma \ref{lemma-compare-with-pullback-flat-proper}.
\item The third holds by Lemma \ref{lemma-proper-flat-base-change}
and the fact that $\mathcal{O}_X = Lf^*\mathcal{O}_Y$.
\item The fourth holds because $\omega_{X/Y}^\bullet = a(\mathcal{O}_Y)$.
\end{enumerate}
Equation (\ref{equation-pre-rigid}) is proved in exactly the same way.
\end{proof}

\end{document}
